\documentclass[11pt]{article}
\usepackage[T1]{fontenc}
\usepackage{lmodern}
\usepackage[margin=1in]{geometry}
\usepackage{amsmath,amssymb,amsthm,mathtools}
\usepackage{bm}
\usepackage{graphicx}
\usepackage{xcolor}
\usepackage{enumitem}
\usepackage{csquotes}
\usepackage[american]{babel}
\usepackage[colorlinks=true,linkcolor=blue,citecolor=blue,urlcolor=blue]{hyperref}
\usepackage{xurl}
\usepackage{booktabs}
\usepackage{microtype}
\emergencystretch=3em

\theoremstyle{plain}
\newtheorem{theorem}{Theorem}[section]
\newtheorem{proposition}[theorem]{Proposition}
\newtheorem{lemma}[theorem]{Lemma}
\newtheorem{corollary}[theorem]{Corollary}
\newtheorem{conjecture}[theorem]{Conjecture}
\newtheorem{observation}[theorem]{Observation}
\theoremstyle{definition}
\newtheorem{definition}[theorem]{Definition}
\newtheorem{question}[theorem]{Question}
\theoremstyle{remark}
\newtheorem{remark}[theorem]{Remark}

\newcommand{\zk}[2]{\zeta_{#1#2}}
\newcommand{\C}{\mathbb{C}}
\newcommand{\R}{\mathbb{R}}
\DeclareMathOperator{\SO}{SO}
\DeclareMathOperator{\Ob}{O}

\title{A verification note on the complex-orthogonal representation of the colored braid groupoid%
\thanks{Comment on arXiv:2606.20473 (I.~E.~Rohozhkin, \emph{Invariants of the Colored Braid Groupoid}).}}
\author{Logan M. Dixon\thanks{Independent Researcher. Email:
\texttt{lmdixon23@gmail.com}. ORCID: 0009-0001-0592-462X.}}
\date{June 2026}

\begin{document}
\maketitle

\begin{abstract}
We record an independent, fully reproducible computational study of the complex-orthogonal
representation $f'$ of the colored braid groupoid introduced by Rohozhkin \cite{rohozhkin2026}
(arXiv:2606.20473). The orthogonality of the flip operators is \emph{structural}:
$\cos^2\phi_{ijkl}+\sin^2\phi_{ijkl}=1$ is a Ptolemy/Pl\"ucker identity holding for every ordering of
the points, so each operator satisfies $A^{\mathsf T}A=I$ over $\C$, with genuine real orthogonality
only when the relevant ratios are positive. The quadrilateral sign rule is \emph{essentially unique}: of the $4^{6}$ assignments exactly two satisfy the pentagon relation for all $120$ orderings (the
paper's $B/C$ rule and its orientation mirror). We also correct a typographical error in the main-text
pentagon equation, given correctly in the appendix, and an erratum in the two displayed worked-example
matrices, which are unit-diagonal but $I+(\text{skew})$ with determinants $\approx 617.52$ and $49$,
hence not orthogonal, although the products of the printed factors are correct. The rational operators
preserve an explicit diagonal form $Q(\Delta_{pqr})=(\zeta_{pq}\zeta_{pr}\zeta_{qr})^{-1}$, so
$f'=DfD^{-1}$ \emph{per flip} with $D=\operatorname{diag}(Q^{1/2})$, giving the unit-diagonal and
determinant-$1$ properties on pure braids; whether the global $B/C$ sign rule realizes this $Q^{1/2}$
branch around every closed loop, so that the transfer applies to the paper's \emph{own} signed $f'$, is
left as a conjecture. For the inner pure braid group at $n=3,4$ the \emph{rational} invariant is exactly
the abelianization $f(\beta)=I+\sum\ell_{ij}(\beta)N_{ij}$ in pairwise linking numbers, so
$\ker f=[PB,PB]$ there, in exact arithmetic on flip sequences that an exact combinatorial check
certifies as legal and resolution-robust; general $n$, the full $PB_{n+3}$, and the signed-$f'$ analogue
remain open. This pins down the incompleteness the paper observes (the Borromean braid is Brunnian,
hence in the kernel). All claims are reproduced by short scripts provided alongside this note; the finite
ones are verified in exact arithmetic.
\end{abstract}

\medskip
\noindent\textbf{2020 Mathematics Subject Classification.} 20F36 (primary); 57K10, 20F14, 13F60, 15A63.

\smallskip
\noindent\textbf{Keywords.} colored braid groupoid; pure braid group; pentagon relation; Ptolemy/Pl\"ucker
identity; complex-orthogonal representation; diagonal similarity; abelianization; linking number;
Brunnian braid; Delaunay flip.

\section{Scope}
This is a verification-and-corrections note, not a new construction. We take the operators of
\cite{rohozhkin2026} (representations of the colored braid groupoid \cite{licata-vertesi}) exactly as
defined there and study them by exact and interval-free
floating-point computation with assertions; every numerical statement below is reproduced by a
named script (\S\ref{sec:repro}). Nothing here bears on the correctness of the paper's main
theorems, which concern the existence of the representations; we sharpen and, in one place,
correct the surrounding apparatus, and we record a structural fact and a uniqueness statement that
the paper does not make explicit.

\section{Setup}
Following \cite{rohozhkin2026}, to each point one attaches a value $\zeta_i$ (the values pairwise
distinct), and to a flip of the diagonal $ik$ to $jl$ inside a quadrilateral with vertices
$i,j,k,l$ one attaches the linear map
\begin{equation}\label{eq:op}
\Delta_{ijk}\mapsto \cos\phi_{ijkl}\,\Delta_{ijl}-\sin\phi_{ijkl}\,\Delta_{jkl},\qquad
\Delta_{ikl}\mapsto \sin\phi_{ijkl}\,\Delta_{ijl}+\cos\phi_{ijkl}\,\Delta_{jkl},
\end{equation}
fixing every other triangle, where with $\zk ab:=\zeta_a-\zeta_b$,
\[
\cos\phi_{ijkl}=(-1)^{s}\sqrt{\frac{\zk il\,\zk jk}{\zk ik\,\zk jl}},\qquad
\sin\phi_{ijkl}=(-1)^{r}\sqrt{\frac{\zk ij\,\zk kl}{\zk ik\,\zk jl}},
\]
the signs $s,r$ depending on the planar type of the quadrilateral (the six types $A$--$F$, with
$s=1$ for type $B$ and $r=1$ for type $C$). The index assignment $i,j,k,l$ is the one of the
paper: $i$ is the vertex of smallest value, $k$ its diagonal opposite, and $j,l$ the remaining
two with $\zeta_j<\zeta_l$.

The paper also gives a \emph{rational} representation $f$ over $\mathbb{Q}$ (the pure-braid invariant
of \cite{rohozhkin-jktr}), with the same support
but rational coefficients,
\begin{equation}\label{eq:op-rational}
\Delta_{ijk}\mapsto \tfrac{\zk il}{\zk ik}\,\Delta_{ijl}+\tfrac{\zk lk}{\zk ik}\,\Delta_{jkl},\qquad
\Delta_{ikl}\mapsto \tfrac{\zk ij}{\zk ik}\,\Delta_{ijl}+\tfrac{\zk jk}{\zk ik}\,\Delta_{jkl},
\end{equation}
to which the complex-orthogonal $f'$ of \eqref{eq:op} is closely related; \S\ref{sec:braid} concerns
both.

\section{Orthogonality is structural (and complex)}
\begin{lemma}[Ptolemy/Pl\"ucker]\label{lem:ptolemy}
For all values $\zeta_i,\zeta_j,\zeta_k,\zeta_l$,
\[
\zk il\,\zk jk+\zk ij\,\zk kl=\zk ik\,\zk jl .
\]
\end{lemma}
\begin{proof}
Expand: $\zk il\zk jk+\zk ij\zk kl=(\zeta_i-\zeta_l)(\zeta_j-\zeta_k)+(\zeta_i-\zeta_j)(\zeta_k-\zeta_l)
=(\zeta_i-\zeta_k)(\zeta_j-\zeta_l)=\zk ik\,\zk jl$ in one line; the cross terms in $\zeta_i\zeta_k$
cancel. \qed
\end{proof}

\begin{remark}
Lemma~\ref{lem:ptolemy} is the classical Pl\"ucker/Ptolemy three-term relation that underlies the
pentagon (Korepanov) literature for operators of this type \cite{korepanov2013}; we record it only to
make the complex-orthogonality below explicit, not as a new identity.
\end{remark}

\begin{proposition}\label{prop:ortho}
For every assignment of pairwise-distinct values, the $2\times2$ active block
$\bigl[\begin{smallmatrix}\cos\phi&\sin\phi\\-\sin\phi&\cos\phi\end{smallmatrix}\bigr]$ of
\eqref{eq:op} satisfies $\cos^2\phi_{ijkl}+\sin^2\phi_{ijkl}=1$, so the flip operator is
complex-orthogonal, $A^{\mathsf T}A=I$. The entries are real (a genuine rotation) iff
$\zk il\zk jk/(\zk ik\zk jl)$ and $\zk ij\zk kl/(\zk ik\zk jl)$ are positive; otherwise the operator
is complex-orthogonal with imaginary entries.
\end{proposition}
\begin{proof}
$\cos^2\phi+\sin^2\phi=\dfrac{\zk il\zk jk+\zk ij\zk kl}{\zk ik\zk jl}=1$ by
Lemma~\ref{lem:ptolemy}. The block is then orthogonal over $\C$, and products of
complex-orthogonal matrices are complex-orthogonal, so the whole image of $f'$ lies in
$\SO(2n+1,\C)$ (each active block has determinant $\cos^2\phi+\sin^2\phi=1$, so the image lies in the
\emph{special} complex-orthogonal group). The reality criterion is immediate from the radicands. \qed
\end{proof}

\begin{remark}
The appendix worked example of \cite{rohozhkin2026} already exhibits imaginary entries
(e.g. $\pm 0.354\,i$); Proposition~\ref{prop:ortho} explains them and locates the image in the
\emph{complex} orthogonal group. This distinction matters for any later construction that would
take advantage of orthogonality (for instance a unitary structure or a trace), and we suggest
stating ``complex-orthogonal'' explicitly.
\end{remark}

\section{Erratum: the main-text pentagon equation}
The main-text pentagon equation (eq.~(2.8) of \cite{rohozhkin2026}) prints its
right-hand side with angles $\phi_{1235},\phi_{2345},\phi_{2345}$: the four-subset $\{1,2,4,5\}$ is
absent and $\{2,3,4,5\}$ appears twice. A pentagon relation uses each of the five four-subsets
$\{1234\},\{1235\},\{1245\},\{1345\},\{2345\}$ exactly once, and the paper's appendix writes the
third factor correctly with $\phi_{1245}$.

\begin{proposition}\label{prop:typo}
With $\zeta=(5,4,3,2,1)$ (the decreasing order assumed in the cited equation), the printed
right-hand side differs from the left-hand side by $\approx 0.26$ in Frobenius norm, while the
appendix form (third factor $\phi_{1245}$) agrees to machine precision.
\end{proposition}
This is a typographical error in the displayed equation only: the appendix uses the corrected
factor $\phi_{1245}$, and we found no other occurrence relying on the erroneous form. Verified
by \texttt{code/pentagon\_typo.py}.

\section{Uniqueness of the sign correction}
The paper notes that the pentagon equation ``holds for the decreasing order and some others, but
not every order,'' which is why the type-dependent signs $s,r$ are introduced. We make this
precise.

\begin{proposition}\label{prop:signs}
Encode a quadrilateral type by the counter-clockwise sequence of value-ranks read from the
smallest-value vertex. Then:
\begin{enumerate}
\item[(a)] With no sign correction, the $5$-cycle (pentagon) relation holds for exactly $60$ of
the $120$ orderings of $\{\zeta_1,\dots,\zeta_5\}$.
\item[(b)] Among all $4^{6}$ assignments of a $(\cos\text{-sign},\sin\text{-sign})\in\{\pm1\}^2$ to
the six types, exactly two make the relation hold for all $120$ orderings. Each has the form
``one type gets a $\cos$-flip, one (other) type gets a $\sin$-flip.'' One is the paper's rule
($\cos$-flip on type $B$, $\sin$-flip on type $C$); the other is its orientation mirror (the two
necklaces are cyclic reverses).
\end{enumerate}
\end{proposition}
Thus the paper's sign rule is not merely \emph{a} correction but is, up to the choice of planar
orientation, the unique one of its form. The count is certified exactly: \texttt{exact\_checks.py}
verifies both winners over all $120$ orderings and exhibits an exact failure witness for each of the
other $4094$ assignments (floating point is used only to locate a candidate witness; the verdict is
exact). Here the $120$ orderings are the assignments of $\{\zeta_1,\dots,\zeta_5\}$ to the five vertices
of one fixed generic convex pentagon; the type-incidences are geometry-aware (each of the six types
arises), but configuration-independence is not separately swept. The $60/120$ and $120/120$ counts are reproduced by \texttt{code/relations.py}.

Explicitly, the six counter-clockwise rank-necklaces correspond to the paper's types $A$--$F$
(Fig.~10) as $A=(1,2,3,4)$, $B=(1,3,2,4)$, $C=(1,2,4,3)$, $D=(1,4,3,2)$, $E=(1,4,2,3)$,
$F=(1,3,4,2)$, and the paper's rule is the $\cos$-flip on $B$ and $\sin$-flip on $C$. The second
solution, the $\cos$-flip on $E=(1,4,2,3)$ and $\sin$-flip on $F=(1,3,4,2)$, is the
orientation-reversed image ($B\leftrightarrow E$, $C\leftrightarrow F$ are cyclic reverses); it is a
genuine alternative under the mirrored planar embedding, not the same rule, which a mirror-orientation
negative test in \texttt{sign\_uniqueness.py} distinguishes.

\section{The groupoid relations, verified}
\begin{proposition}\label{prop:relations}
For the operators \eqref{eq:op} with the sign rule of Proposition~\ref{prop:signs}, all four
defining relations of $\mathcal G^{4}_{n+3}$ \cite{manturov-nikonov,fedoseev-manturov-nikonov} hold: the identity and $4$-cycle relations by
construction; the $5$-cycle relation for all $120$ orderings (Proposition~\ref{prop:signs}); and
far-commutativity, namely that flips in quadrilaterals with disjoint interiors commute. This is checked
numerically over all $8!$ orderings of an eight-point configuration (residual $0$), with an exact
representative check and a structural disjoint-support argument (\texttt{exact\_checks.py}).
\end{proposition}
Verified by \texttt{code/relations.py} and, exactly, by \texttt{code/exact\_checks.py}.
Far-commutativity is in fact a proposition, not merely a tested instance: if two flips occur in
quadrilaterals with disjoint vertex sets, their four affected triangles $\{ijk,ikl,ijl,jkl\}$ are
disjoint from the other flip's four, so the two operators act on complementary coordinate blocks of
$\mathbb{C}^{(\mathbf T)}$ and therefore commute; the $8!$-ordering check confirms one instance.

\section{Pure-braid images and the trace obstruction (computational)}\label{sec:braid}
We reconstruct the paper's worked pure-braid example flip-by-flip. The appendix presents
$f'(b_{\zeta_4\zeta_7})$ ($n=4$, $\zeta_l=l$) as a product of sixteen named flip matrices
$A_{ikjl}$ and prints two of them explicitly. Reading the subscript $A_{ikjl}$ as the operator
indices positionally and taking the \emph{principal} complex square root (with no separate
$(-1)^s,(-1)^r$ sign flip, the branch being automatic) reproduces both printed matrices to eight
digits ($A_{4615}$: $\cos=0.79056942$, $\sin=0.61237244$; $A_{1524}$: $\cos=1.06066017$,
$\sin=0.35355339\,i$). This locks the per-flip convention.

Composing the sixteen factors (rightmost first) reproduces the paper's stated global property of
pure-braid matrices (determinant $1$ and unit diagonal) to $\sim 10^{-16}$, \emph{for the
rational representation} $f$ of \eqref{eq:op-rational}; the reversed composition order and the
reconstruction of the initial triangulation are validated jointly by this match against the paper's
own assertion.

\begin{remark}[Why the na\"ive trace fails]
\cite{rohozhkin2026} states that \emph{every} pure-braid matrix has unit diagonal; granting this,
$\operatorname{tr} f(\beta)=2n+1$ for all pure $\beta$, so the trace carries no information and is
not a candidate knot/link invariant. We do not prove the universal statement here; we confirm it on
the worked example $b_{\zeta_4\zeta_7}$ (where $\operatorname{tr} f=9$), which is the obstruction
noted in \cite{rohozhkin2026}.
It motivates deforming the operators so that the diagonal, determinant, or trace can vary.
\end{remark}

Under naive matrix composition the \emph{complex-orthogonal} product $f'(b_{\zeta_4\zeta_7})$ is
\emph{not} unit-diagonal ($\max_i|M_{ii}-1|=3.0$), although \cite{rohozhkin2026} states the property
for both representations. Section~\ref{sec:similar} resolves this for the $Q$-orthonormalized
conjugate: per flip $f'=DfD^{-1}$ with $D=\operatorname{diag}(Q^{1/2})$, so the $Q$-orthonormalized
$f'(\beta)$ has the same diagonal as $f(\beta)$, the naive product failing only because the square-root
branch defining $D$ was not tracked consistently around the loop. Whether this $Q$-orthonormalized
operator coincides with the paper's $B/C$-signed $f'(\beta)$ around the loop is the branch conjecture of
Proposition~\ref{prop:similar}'s remark. The braid-simulation pipeline of Section~\ref{sec:kernel} (which moves the inner points, reads off
the Delaunay flips, and composes the operators) reproduces the braid relations, far-commutativity,
and the Borromean collapse $(\sigma_1\sigma_2^{-1})^3\mapsto I$ in these diagnostic checks, and is what
Section~\ref{sec:kernel} uses to \emph{locate} the flip sequences on the inner pure braid group.

\section{The two representations are diagonally similar}\label{sec:similar}
The discrepancy above has a structural explanation: the rational and complex-orthogonal representations
differ by a diagonal change of basis.

\begin{proposition}\label{prop:similar}
Let $Q$ be the diagonal form on $\mathbb{C}^{(\mathbf T)}$ with
$Q(\Delta_{pqr})=(\zk pq\,\zk pr\,\zk qr)^{-1}$, the reciprocal product of the triangle's three
pairwise differences. Every rational flip operator \eqref{eq:op-rational} is an isometry of $Q$:
$f^{\mathsf T}Q_p\,f=Q_t$ between the source and target triangulations. Equivalently, writing
$D=\operatorname{diag}\big(Q^{1/2}\big)$, the complex-orthogonal operator \eqref{eq:op} is
$f'=D\,f\,D^{-1}$.
\end{proposition}
\begin{proof}
A flip changes only the four triangles $\Delta_{ijk},\Delta_{ikl}\to\Delta_{ijl},\Delta_{jkl}$;
every other triangle is fixed and contributes equal diagonal entries to $Q_t$ and $Q_p$. On the
$2\times2$ active block, $f^{\mathsf T}\operatorname{diag}(Q_{ijl},Q_{jkl})\,f
=\operatorname{diag}(Q_{ijk},Q_{ikl})$ is a rational identity in $\zeta_i,\zeta_j,\zeta_k,\zeta_l$,
verified by symbolic computation (\texttt{code/diagonal\_form.py}, residual $0$). Dividing the
$(\mathrm{row},\mathrm{col})$ entry of $f$ by $\sqrt{Q(\mathrm{col})/Q(\mathrm{row})}$ turns
$\zk il/\zk ik$ into $\sqrt{\zk il\,\zk jk/(\zk ik\,\zk jl)}=\cos\phi_{ijkl}$ and likewise for the
other three entries, which is precisely \eqref{eq:op}; hence $f'=DfD^{-1}$. \qed
\end{proof}

\begin{corollary}\label{cor:unitdiag}
For every flip word $\beta$, $f(\beta)$ is an isometry of $Q$. A composite of $Q$-isometries is again
one: for two consecutive flips the triangle they share carries a single value of $Q$, so the first
flip's target form is the second's source form and the two cancel; Proposition~\ref{prop:similar} is
the one-flip case, and the two-flip case is checked symbolically
(\texttt{code/diagonal\_form.py}). For a pure braid the start and end triangulations coincide,
so $f(\beta)$ preserves $Q_{t_0}$, and with $D=\operatorname{diag}(Q_{t_0}^{1/2})$ the conjugate
$f'(\beta)=D\,f(\beta)\,D^{-1}$ is complex-orthogonal; being a diagonal conjugate of $f(\beta)$, it has
the same main diagonal and the same determinant. Granting the property of \cite{rohozhkin2026} that
$f(\beta)$ has unit diagonal and determinant $1$ (Section~\ref{sec:braid}, confirmed on
$b_{\zeta_4\zeta_7}$), so does $f'(\beta)$ , where $f'(\beta)$ here denotes the $Q$-orthonormalized
conjugate $D\,f(\beta)\,D^{-1}$, whose agreement with the paper's $B/C$-signed operator around the loop
is the branch identification conjectured in the remark below. The unit-diagonal and determinant-$1$
transfer is thus certified for the $Q$-branch conjugate and conditional, through that conjecture, for the
paper's \emph{own} $f'(\beta)$. This is the property naive composition appeared to violate.
\end{corollary}

\begin{remark}
The factor $D$ involves $Q^{1/2}$, a square root with a sign per triangle (the form $Q$ is
antisymmetric under relabeling a triangle's vertices). Composing cos/sin matrices with independently
chosen principal branches need not keep $D$ single-valued around a closed loop. We therefore
\emph{conjecture} that the paper's $B/C$ sign rule is exactly the branch bookkeeping for $Q^{1/2}$;
this is a structural inference, not proved here at the level of the algebraic identity above. What is
checked: on $b_{\zeta_4\zeta_7}$, numerically recovering the form $Q$ that the loop preserves and
conjugating yields a matrix that is, to machine precision, complex-orthogonal, unit-diagonal, and of
determinant $1$ (\texttt{diagonal\_form.py}); a confirmation on the example, not a proof of the
global branch identification.
\end{remark}

\begin{remark}[an erratum in the worked examples]
The paper prints, for $\zeta_l=l$, the full $9\times9$ matrices $f'(b_{\zeta_4\zeta_7})$ and
$f'(b_{\zeta_5\zeta_6})$ as products of named flip operators followed by a displayed result. The
displayed result matrices are incorrect: each is unit-diagonal but is $I+(\text{skew-symmetric})$,
hence \emph{not} complex-orthogonal and with determinant $\approx 617.52$ and $49$ respectively, contradicting
the paper's own properties (orthogonal; determinant $1$). The paper's \emph{factors} are fine:
multiplying them in printed order gives, in both cases, a complex-orthogonal, unit-diagonal,
determinant-$1$, rank-$2$ one-step unipotent matrix, consistent with the rational computation via
Proposition~\ref{prop:similar} (\texttt{appendix\_full\_matrix.py}). Their rank-$2$ unipotent form also
corroborates Observation~\ref{obs:abelian} on the inner-outer generators $A_{47},A_{56}$.
\end{remark}

\begin{question}
What is the kernel of $f'$ on $\mathrm{ColB}(n)$? Granting the branch conjecture above it equals the
kernel of $f$, whose image lies in $\SO(2n+1,\C)$ and is unit-diagonal on pure braids; characterizing
the subgroup generated, and the kernel, would explain any collapse of the invariant structurally
rather than by a single computation.
\end{question}

\section{The kernel: the invariant is the abelianization}\label{sec:kernel}
The question above has, at least for small cases, a complete answer. Realizing braids as
motions of the inner points and recording the Delaunay flips \cite{aurenhammer} (\texttt{code/kernel\_probe.py};
the realization must braid strands by \emph{position}, after which trivially trivial braids, the
braid relation $\sigma_1\sigma_2\sigma_1=\sigma_2\sigma_1\sigma_2$, far commutation, and the
Borromean braid $(\sigma_1\sigma_2^{-1})^3$ all map to the identity), we compute the representation on
the pure braid group of the inner strands. The exact computation (\texttt{code/exact\_kernel.py}) applies the \emph{rational} operators $f$ to the
flip sequence the Delaunay pipeline records, so the structure below is established for $f$ in exact
arithmetic, conditional on that float-located sequence; the floating-point pipeline
(\texttt{code/float\_kernel.py}) applies the complex operators and agrees diagnostically. The transfer to
the signed $f'$ is conditional on the diagonal-similarity conjecture (Proposition~\ref{prop:similar}).
For $n=3$ and $n=4$ the structure is uniform:

\begin{observation}[verified in exact arithmetic for the inner $PB_n$, $n=3,4$, given the Delaunay-located flip sequences]\label{obs:abelian}
Each pure-braid generator $A_{ij}$ (linking of strands $i,j$) maps to $I+N_{ij}$ where
$N_{ij}^2=0$, $\operatorname{rank}N_{ij}=2$, \emph{all products} $N_{ij}N_{kl}=0$, and the $N_{ij}$
are linearly independent and mutually commuting. Hence $f$ is purely additive,
\[
   f(\beta)=I+\sum_{i<j}\ell_{ij}(\beta)\,N_{ij},
\]
with $\ell_{ij}(\beta)$ the exponent sum of $A_{ij}$ in $\beta$, i.e.\ the pairwise linking number.
The number of independent $N_{ij}$ equals $\binom{n}{2}=\dim PB_n^{\mathrm{ab}}$, so $f$ is the
abelianization $PB_n\to\mathbb{Z}^{\binom n2}$ on the tested group, whence
\[
   \ker f=[PB_n,PB_n]\qquad(\text{inner }PB_n,\ n=3,4),
\]
in exact arithmetic on the Delaunay-located flip sequences; $\ker f'=[PB_n,PB_n]$ holds conditionally on
the diagonal-similarity conjecture (Proposition~\ref{prop:similar}).
\end{observation}

Thus the invariant records exactly the pairwise linking numbers and is blind to the entire commutator
subgroup. This makes precise the incompleteness \cite{rohozhkin2026} observes: the Borromean braid is
Brunnian \cite{kassel-turaev} (so all pairwise linking numbers vanish), hence lies in $[PB,PB]$, hence maps to $I$. The constant trace of
Section~\ref{sec:braid} is the same fact, since $I+\sum\ell_{ij}N_{ij}$ has trace $2n+1$ regardless of
the $\ell_{ij}$. Any trace- or determinant-based knot invariant must therefore deform the construction
away from this abelian, unipotent image.

This is verified in exact arithmetic (\texttt{exact\_kernel.py}, sympy: the braid's combinatorial flip
sequence with exact rational operators) for the inner pure braid group at $n=3,4$; \texttt{float\_kernel.py}
is the floating-point version. Generators that loop an inner strand around a fixed point are also rank-2
one-step unipotents, consistent with the same description for the full $PB_{n+3}$, but we neither
realize them with the present tool nor prove the statement for general $n$.

\section{Verification and reproducibility}\label{sec:repro}

\paragraph{Artifact of record.} The accompanying repository
(\url{https://github.com/lmdixon23/cbg-2026-representation-checks}) ships the code; its
\texttt{verification/STATUS-LEDGER.md} records the evidence type and status of every claim above, and
\texttt{run\_all.sh} runs the suite in claim order.

\paragraph{Reproduction protocol.} The computational claims can be reproduced from the repository with
\texttt{bash run\_all.sh} (Python, \texttt{numpy}/\texttt{scipy}/\texttt{sympy}); each script is
self-contained and prints its checks. The exact certificates use no floating point and no tolerance.

\paragraph{Claim-to-artifact map.} \texttt{code/operators.py} (operators, Lemma~\ref{lem:ptolemy});
\texttt{code/relations.py} and \texttt{code/exact\_checks.py} (the four relations; the $60/120$ and
$120/120$ counts; sign uniqueness); \texttt{code/pentagon\_typo.py} (Proposition~\ref{prop:typo});
\texttt{code/sign\_uniqueness.py} (Proposition~\ref{prop:signs});
\texttt{code/appendix\_reconstruction.py} (the convention lock and the rational pure-braid property of
\S\ref{sec:braid}); \texttt{code/diagonal\_form.py} (Proposition~\ref{prop:similar});
\texttt{code/appendix\_full\_matrix.py} (the appendix worked-example erratum); \texttt{code/float\_kernel.py}
and \texttt{code/exact\_kernel.py} (Observation~\ref{obs:abelian}, floating point and exact, inner $PB_n$,
$n=3,4$); \texttt{code/check\_flip\_sequence.py} (the exact combinatorial flip-sequence certificate).
\texttt{code/kernel\_probe.py} is a diagnostic only (braid pipeline; Borromean $\to I$) and supports no
property claim. Each load-bearing claim carries an evidence type in the ledger (algebraic identity, finite
exhaustive, checker-verified, or strategic inference).

\paragraph{Independent checks.} \texttt{BLIND\_CHECK.py} recomputes the headline quantities from the
formulas alone, importing nothing from \texttt{code/}, so a transcription or convention error would surface
as a disagreement. The load-bearing identities were additionally re-derived by hand (the Ptolemy lemma
(Lemma~\ref{lem:ptolemy}), the single-flip form identity $f^{\mathsf T}Qf=Q$ behind
Proposition~\ref{prop:similar}, and the two printed flip matrices) and an independent pass re-ran
the suite in a separate environment and carried out the literature checks.

\paragraph{Limitations and open obligations.} The finite claims are exact; the conjectural items rest on no
machine assertion. The identification of the paper's $B/C$ sign rule with the $Q^{1/2}$ branch, the
global transfer to the paper's own signed $f'$, is an explicit conjecture (Proposition~\ref{prop:similar}
and its remark). The kernel result is exact for the rational $f$ on the inner $PB_n$, $n=3,4$, conditional
on the Delaunay-located flip sequence (verified legal and resolution-robust, but not yet derived
geometry-free); the statement for general $n$, the full $PB_{n+3}$, and the signed $\ker f'$ remain open.

\paragraph{AI assistance.} AI systems were used as research tools during this project, including Claude by
Anthropic and GPT by OpenAI. They assisted with candidate generation, code drafting, adversarial review, and
verification-ledger preparation. Their outputs were not treated as authoritative: candidate constructions,
code, constants, and literature summaries were accepted only after source checking, independent
recomputation, author-executed verification runs, or hand derivation where appropriate. The author set the
research direction and acceptance criteria, ran the verification pipeline, and bears sole responsibility for
all claims.

\section*{Declarations}
\paragraph{Funding.} No funding was received for this work.
\paragraph{Competing interests.} The author declares no competing interests.
\paragraph{Data and code availability.} The code, verification scripts, and status ledger are available in
the accompanying repository (\url{https://github.com/lmdixon23/cbg-2026-representation-checks}); an archival
DOI will be added at submission.

\begin{thebibliography}{9}
\bibitem{rohozhkin2026} I.~E.~Rohozhkin, \emph{Invariants of the Colored Braid Groupoid},
preprint, 2026, \href{https://arxiv.org/abs/2606.20473}{arXiv:2606.20473}.
\bibitem{korepanov2013} I.~G.~Korepanov and N.~M.~Sadykov, \emph{Pentagon relations in direct sums
and Grassmann algebras}, SIGMA Symmetry Integrability Geom. Methods Appl. \textbf{9} (2013), Paper 030,
16 pp. \href{https://doi.org/10.3842/SIGMA.2013.030}{doi:10.3842/SIGMA.2013.030}; \href{https://arxiv.org/abs/1212.4462}{arXiv:1212.4462}.
\bibitem{rohozhkin-jktr} I.~E.~Rohozhkin, \emph{Pentagon equations, Delaunay triangulations and the
pure braid group invariant}, J. Knot Theory Ramifications (2025).
\bibitem{licata-vertesi} J.~Licata and V.~V\'ertesi, \emph{Liftable braids and colored braid groupoid},
preprint, 2025, \href{https://arxiv.org/abs/2508.05146}{arXiv:2508.05146}.
\bibitem{manturov-nikonov} V.~O.~Manturov and I.~M.~Nikonov, \emph{The groups $\Gamma^4_n$, braids, and
3-manifolds}, preprint, 2023, \href{https://arxiv.org/abs/2305.06316}{arXiv:2305.06316}.
\bibitem{fedoseev-manturov-nikonov} D.~A.~Fedoseev, V.~O.~Manturov, and I.~M.~Nikonov, \emph{Manifolds
of triangulations, braid groups of manifolds, and the groups $\Gamma^k_n$}, preprint, 2020,
\href{https://arxiv.org/abs/1912.02695}{arXiv:1912.02695}.
\bibitem{kassel-turaev} C.~Kassel and V.~Turaev, \emph{Braid Groups}, Graduate Texts in Mathematics
\textbf{247}, Springer, 2008.
\bibitem{aurenhammer} F.~Aurenhammer, R.~Klein, and D.-T.~Lee, \emph{Voronoi Diagrams and Delaunay
Triangulations}, World Scientific, 2013.
\end{thebibliography}

\end{document}
